% Lecture Template for ENGR 1120 030 031 E01- Tristan Hill - Spring 2017 - Summer 2018
% 
% Introduction to MATLAB 
%
% THe Command WIndow

% Document settings
\documentclass[11pt]{article}
\usepackage[margin=1in]{geometry}
\usepackage[pdftex]{graphicx}
\usepackage{multirow}
\usepackage{setspace}
\usepackage{hyperref}
\usepackage{color,soul}
\usepackage{fancyvrb}
\usepackage{framed}
\usepackage{wasysym}
\usepackage{multicol}

\pagestyle{plain}
\setlength\parindent{0pt}
\hypersetup{
    bookmarks=true,         % show bookmarks bar?
    unicode=false,          % non-Latin characters in Acrobat’s bookmarks
    pdftoolbar=true,        % show Acrobat’s toolbar?
    pdfmenubar=true,        % show Acrobat’s menu?
    pdffitwindow=false,     % window fit to page when opened
    pdfstartview={FitH},    % fits the width of the page to the window
    pdftitle={My title},    % title
    pdfauthor={Author},     % author
    pdfsubject={Subject},   % subject of the document
    pdfcreator={Creator},   % creator of the document
    pdfproducer={Producer}, % producer of the document
    pdfkeywords={keyword1} {key2} {key3}, % list of keywords
    pdfnewwindow=true,      % links in new window
    colorlinks=true,       % false: boxed links; true: colored links
    linkcolor=red,          % color of internal links (change box color with linkbordercolor)
    citecolor=green,        % color of links to bibliography
    filecolor=magenta,      % color of file links
    urlcolor=blue           % color of external links
}

\definecolor{mygreen}{rgb}{0, .39, 0}

%\definecolor{dred}{#8B0000}
% [153,50,204] - dark orchid
\definecolor{mypurple}{rgb}{0.6,0.1961,0.8}
%[139,69,19] - saddle brown
\definecolor{mybrown}{rgb}{0.5451,0.2706,0.0745}

\definecolor{mygray}{rgb}{.6, .6, .6}

\setulcolor{red} 
\setstcolor{green} 
\sethlcolor{mygray} 

\newcommand{\VA}{\vspace{2mm}}
\newcommand{\VB}{\vspace{5mm}} 
 
\newcommand{\R}{\color{red}}
\newcommand{\B}{\color{blue}}
\newcommand{\K}{\color{black}}
\newcommand{\G}{\color{mygreen}}
\newcommand{\PR}{\color{mypurple}}

% assignment number 
\newcommand{\NUM}{1 } 
\newcommand{\VSpaceSize}{2mm} 
\newcommand{\HSpaceSize}{2mm} 

\begin{document}

\textbf{ \LARGE ENGR 1120 Lecture Chapter \NUM \\\\- More Numerical Expressions} \\
\Large
\begin{itemize}

	\item \textbf{ \LARGE Useful Commands( type in Command Window)}\\

		
		\scalebox{1.5}{{\fontfamily{qcr}\selectfont  \hspace{5mm} >>  clear variables}} \\

		
		\scalebox{1.5}{{\fontfamily{qcr}\selectfont  \hspace{5mm} >>  clc}} \\\\

	\item \textbf{ \LARGE 1.3 - Variables and Assignment} \\
		\Large
		\begin{itemize}
			\item \Large{A variable is a container for storing values in the RAM. The value that is stored has a {\it type} (1.3.3). We will begin with {\it floating point} values.} \vspace{50mm}
			
		\scalebox{2}{$variable = expression$}	\vspace{50mm}

		
		\newpage
			\item \LARGE{Basic Use of a \PR{\bf Variable}} \\
				\begin{itemize}
					
					\item {\bf Assign} a Value  \vspace{35mm}

					\item {\bf Access} the Value  \vspace{35mm}

					\item {\bf Re-Assign} a Value \vspace{35mm}
					
					\item Common Errors \vspace{35mm}
				

				\end{itemize}


		\end{itemize}
	\newpage
	\item \textbf{ \LARGE 1.4 - Numerical Expressions}
		\begin{itemize}
			\item \Large{We will begin with numerical expressions. This is how we do typical math computations in MATLAB.} \vspace{10mm}

			\item To do this we need to learn about {\it operator precedence}. \\

				%\begin{multicols}{2}				

				\scalebox{2}{$(\hspace{3mm} )$} \hspace{15mm}\Large{Parenthesis} \\\\

				\scalebox{2}{$\wedge$} \hspace{15mm}\Large{Exponents} \\\\

				\scalebox{2}{$-$} \hspace{15mm}\Large{Negation} \\\\

				\scalebox{2}{$*\hspace{3mm} /$} \hspace{15mm}\Large{Multiplication and Division}\\\\				
	
				 \scalebox{2}{$+\hspace{3mm} -$} \hspace{15mm}\Large{Addition and Subtraction}\\\\
				
				%\end{multicols}	
		
				
		\end{itemize}

\newpage
	\item \textbf{ \LARGE 1.4.3 -  Built-In MATLAB functions}\\\\
	 \includegraphics[scale=.15]{lecture2_fig1.png} \hspace{35mm}\includegraphics[scale=.30]{lecture2_fig2.png}\\\\  
		
		\begin{itemize}
			\item \textbf{\Large What is a function?} \vspace{60mm}
			\item \textbf{\Large How do we use a function?} \\\\
			 \scalebox{1.5}{{\fontfamily{qcr}\selectfont  \hspace{1mm} >> [y1,y2,...,ym]=fun(x1,x2,...,xn) }} \vspace{100mm}
			\item \textbf{\Large Typical Mathematics Functions} \\	
				\begin{itemize}
					\item \scalebox{1.5}{{\fontfamily{qcr}\selectfont  \hspace{5mm} sqrt()}} \\
					\item \scalebox{1.5}{{\fontfamily{qcr}\selectfont  \hspace{5mm} exp()}} \\
					\item \scalebox{1.5}{{\fontfamily{qcr}\selectfont  \hspace{5mm} log()}} \\
					\item \scalebox{1.5}{{\fontfamily{qcr}\selectfont  \hspace{5mm} log2()}} \\
					\item \scalebox{1.5}{{\fontfamily{qcr}\selectfont  \hspace{5mm} log10()}} \\
				\end{itemize}	
				\item \textbf{\Large Other Useful Functions} \\	
				\begin{itemize}
					\item \scalebox{1.5}{{\fontfamily{qcr}\selectfont  \hspace{5mm} round()}} \\
					\item \scalebox{1.5}{{\fontfamily{qcr}\selectfont  \hspace{5mm} floor()}} \\
					\item \scalebox{1.5}{{\fontfamily{qcr}\selectfont  \hspace{5mm} int8()}} \\
					\item \scalebox{1.5}{{\fontfamily{qcr}\selectfont  \hspace{5mm} sign()}} \\
					\item \scalebox{1.5}{{\fontfamily{qcr}\selectfont  \hspace{5mm} mod()}} \\
					\item \scalebox{1.5}{{\fontfamily{qcr}\selectfont  \hspace{5mm} rem()}} \\
	
				\end{itemize}	
			\item \textbf{\Large The Built in Help }\\	
				\begin{itemize}
					\item \LARGE{use the help to get information about the built in functions } \\\\
					\item \scalebox{1}{{\fontfamily{qcr}\selectfont  \hspace{5mm} >> help sqrt()}} \\
	
				\end{itemize}	
				\newpage


	
		\end{itemize}






	\newpage
	\item \textbf{ \LARGE 1.4.4 - Constants} \\
		\begin{itemize}
			\item \LARGE{Several useful constants are built into MATLAB.} \\

				\begin{itemize}
					
					\item \scalebox{1.5}{{\fontfamily{qcr}\selectfont  \hspace{5mm} pi}} \\
					\item \scalebox{1.5}{{\fontfamily{qcr}\selectfont  \hspace{5mm} i}} \\
					\item \scalebox{1.5}{{\fontfamily{qcr}\selectfont  \hspace{5mm} j}} \\
					\item \scalebox{1.5}{{\fontfamily{qcr}\selectfont  \hspace{5mm} inf}} \\
					\item \scalebox{1.5}{{\fontfamily{qcr}\selectfont  \hspace{5mm} NaN}} \\\\

				\end{itemize}



		\end{itemize}

	

\item \textbf{ \LARGE 1.4.5 - Random Numbers} \\
		\begin{itemize}
			\item \LARGE{Sometime it is useful generate random data in MATLAB.} \\

				\begin{itemize}
					
					\item \scalebox{1.5}{{\fontfamily{qcr}\selectfont  \hspace{5mm} rand()}} \\
					\item \scalebox{1.5}{{\fontfamily{qcr}\selectfont  \hspace{5mm} randi()}} \\\\\\
		
				\end{itemize}



		\end{itemize}
\item \textbf{ \LARGE 1.5 and 1.6 - We will cover after we cover chapter 2} \\
		\newpage
\item \textbf{ \LARGE Let us solve an example like the lab} \\
\newpage
\item \textbf{ \LARGE example } \\
\end{itemize}


	

\end{document}



